% Informe de avance para el dueño del producto: dónde estamos.
% Se compila con:   cd informes && pdflatex estado-2026-11-04.tex   (dos veces)
% La letra es Times (mathptmx: la Times de URW). Times New Roman como tal no
% está instalada donde se compila; esta es la misma familia de letra.
\documentclass[12pt,a4paper]{article}

\usepackage[T1]{fontenc}
\usepackage[utf8]{inputenc}
% Times (la versión libre de URW, con las mismas formas y medidas)
\usepackage{mathptmx}
\usepackage[spanish,es-noshorthands]{babel}
\usepackage[a4paper,top=2.5cm,bottom=2.5cm,left=2.8cm,right=2.8cm]{geometry}
\usepackage{setspace}
\onehalfspacing
\setlength{\emergencystretch}{3em}   % los nombres de archivo largos no se cortan
\usepackage{microtype}
\usepackage[dvipsnames]{xcolor}
\usepackage{enumitem}
\setlist{itemsep=0.35em, topsep=0.5em}
\usepackage{titlesec}
\titleformat{\section}{\Large\bfseries\color{MidnightBlue}}{\thesection.}{0.6em}{}
\titleformat{\subsection}{\large\bfseries}{}{0em}{}
\titlespacing*{\section}{0pt}{1.6em}{0.6em}
\titlespacing*{\subsection}{0pt}{1.1em}{0.3em}
\usepackage[most]{tcolorbox}
\newtcolorbox{recuadro}[1]{colback=MidnightBlue!4, colframe=MidnightBlue!60, fonttitle=\bfseries,
  title={#1}, arc=2mm, boxrule=0.6pt, left=3mm, right=3mm, top=2mm, bottom=2mm, before skip=1em, after skip=1em}
\newtcolorbox{atencion}[1]{colback=Orange!6, colframe=Orange!75!black, fonttitle=\bfseries,
  title={#1}, arc=2mm, boxrule=0.6pt, left=3mm, right=3mm, top=2mm, bottom=2mm, before skip=1em, after skip=1em}
\usepackage{fancyhdr}
\pagestyle{fancy}
\fancyhf{}
\fancyhead[L]{\small\color{gray}Micros Tempo · Dónde estamos}
\fancyhead[R]{\small\color{gray}4 de noviembre de 2026}
\fancyfoot[C]{\small\thepage}
\renewcommand{\headrulewidth}{0.3pt}
\usepackage[hidelinks]{hyperref}

\begin{document}

\begin{titlepage}
  \centering
  \vspace*{4cm}
  {\Huge\bfseries Dónde estamos\par}
  \vspace{0.8cm}
  {\Large Lo que cambió en el sistema de control de flota\\
  entre el 31 de octubre y el 4 de noviembre de 2026\par}
  \vspace{2.5cm}
  {\large Micros Tempo\par}
  \vfill
  {\normalsize Escrito para leer de corrido, sin términos técnicos.\\
  El detalle de cada cambio, para quien lo programe, está en\\
  el archivo \emph{CAMBIOS-PARA-REVISAR.md} del proyecto.\par}
  \vspace{1cm}
\end{titlepage}

\section{En pocas palabras}

En estos cinco días el sistema dio un salto. Se terminó la lista de cosas
que hacían falta antes de salir con una cooperativa de verdad, y después se
revisó todo lo nuevo con ojos frescos para encontrar lo que pudiera fallar
en la calle. Lo que se encontró se corrigió.

Hoy el sistema está \textbf{listo para el piloto en Juliaca}. Lo que falta
ya no es programar: son pasos tuyos (compilar la aplicación, probarla en un
teléfono, configurar un par de cosas en los paneles de los servicios) y
algunas decisiones que te tocan como dueño del producto.

\begin{recuadro}{Lo más importante de estos días}
\begin{itemize}[leftmargin=1.2em]
  \item Una \textbf{emergencia (SOS)} ya no puede pasar desapercibida: suena
  en Despacho hasta que alguien la atiende, y el chofer ve que lo vieron.
  Si no hay internet, el chofer puede llamar o mandar un mensaje de texto
  a la central.
  \item Los arreglos de la aplicación del chofer \textbf{llegan solos} a los
  teléfonos, sin reinstalar.
  \item Un reinicio del servidor ya \textbf{no le borra la vuelta} a las
  combis que estaban en ruta.
  \item Hay una \textbf{copia de seguridad fuera del servidor}, cifrada, y una
  forma probada de volver a ella.
  \item Quedó preparada la parte de \textbf{datos personales} que pide la ley
  peruana.
  \item Una \textbf{revisión a fondo} encontró y corrigió 27 problemas, cinco
  de ellos graves.
\end{itemize}
\end{recuadro}

\section{Lo nuevo, según quién lo usa}

\subsection{Para el chofer y el cobrador}

\textbf{Saber que Despacho vio su emergencia.} Cuando el chofer pide ayuda y
Despacho la marca como atendida, en el teléfono aparece «Despacho vio tu SOS»,
con quién lo atendió y a qué hora. Antes el chofer no tenía forma de saber si
alguien lo había visto.

\textbf{Pedir ayuda sin internet.} Si el SOS no sale porque el teléfono no
tiene datos, la aplicación lo dice claramente y ofrece dos botones: llamar a
la central o mandarle un mensaje de texto ya escrito, con la combi y un enlace
al mapa de dónde está. También muestra los números de emergencia (policía,
ambulancia, bomberos). El teléfono de la central lo carga la gerencia de cada
cooperativa.

\textbf{Arreglos sin reinstalar.} Hasta ahora, cada mejora de la aplicación
obligaba a repartir un archivo nuevo y que cada chofer lo instalara. Desde la
versión 0.3.0 los arreglos llegan solos. La aplicación los baja en silencio y
los aplica \textbf{sólo cuando el chofer no está trabajando}: nunca en medio de
una vuelta. Hay excepciones (por ejemplo, si se necesita un permiso nuevo del
teléfono); para eso sigue haciendo falta instalar.

\textbf{Pedir una copia de sus datos.} La ley le da derecho a pedir todo lo
que el sistema tiene de él. Despacho ahora lo baja con un botón (ver más
abajo).

\subsection{Para Despacho}

\textbf{El SOS suena.} Una emergencia hace sonar una alarma en el panel, hace
parpadear la pestaña del navegador y muestra un aviso en la computadora,
aunque Despacho esté en otra ventana. Sigue sonando hasta que alguien toca
ATENDIDO. Si hay varios paneles abiertos, todos se enteran de quién lo
atendió. Al empezar el turno, el operador tiene que tocar una vez «Activar
alarma de SOS» y aceptar las notificaciones: los navegadores no dejan sonar
nada sin ese primer toque.

\textbf{La batería del teléfono de cada combi.} En la tarjeta de cada combi
se ve la batería del teléfono del chofer, si está cargando, y si el teléfono
le está aplicando «ahorro de batería» a la aplicación. Esa última es la causa
número uno de «la combi no aparece en el mapa». Antes había que llamar al
chofer para averiguarlo.

\textbf{Corregir horas a mano.} Si un chofer se olvidó de cerrar el turno, o
el teléfono se apagó, Despacho puede corregir la hora de entrada o de salida,
anular un turno equivocado o agregar uno que faltó. Siempre con un motivo
escrito, que queda registrado junto con quién lo hizo. Sólo se pueden tocar
los turnos del último mes: lo anterior ya se liquidó.

\textbf{Cargar a todos de una vez.} En lugar de dar de alta a cada chofer y
cobrador por separado, Despacho puede subir la planilla que ya tenga (de
Excel, de Google Sheets, de LibreOffice, o simplemente copiar y pegar las
celdas). El sistema primero muestra qué va a hacer con cada fila, sin guardar
nada; si todo está bien, se confirma y se crean las cuentas, con sus claves
para entregar. Hay una planilla de ejemplo para bajar.

\textbf{Guardar algo como evidencia.} Las fotos y las notas de voz del chat
se conservan 30 días. Si una foto o un audio sirve de prueba de algo (un
choque, un pasajero agresivo), Despacho lo marca como evidencia, escribe por
qué, y se conserva un año. Los demás paneles ven la marca. Los choferes no
ven el motivo.

\textbf{La Actividad, ordenada.} El registro de «quién hizo qué» ahora se
puede filtrar en tres partes: la \emph{gestión} (altas, bajas, cambios de
clave, correcciones), la \emph{operación} (desvíos, emergencias, cambios de
combi) y los \emph{accesos} (quién entró al sistema). Antes los ingresos
llenaban el registro en días y se llevaban puestas las altas y las bajas, que
son justo lo que hay que poder consultar meses después. Ahora cada cosa se
guarda lo que corresponde: los accesos tres meses, la operación un año, la
gestión dos años.

\textbf{El consentimiento de cada chofer.} En Personas, Despacho ve cuántos
choferes y cobradores firmaron el formulario de consentimiento de datos
personales y cuántos faltan. Anota la fecha del papel firmado, y si alguno lo
retira, también queda anotado. El formulario se imprime desde el mismo
sistema.

\textbf{Bajar todos los datos de una persona.} En Personas, al editar a
alguien, hay un botón que baja un archivo con todo lo que el sistema tiene de
esa persona: su cuenta, sus turnos, sus mensajes con sus fotos y audios, sus
vueltas y desvíos (sólo de cuando manejaba él, no de la combi con otro
chofer) y lo que hizo en el sistema. Trae una hoja que explica cada parte. Es
lo que se le entrega si pide sus datos.

\subsection{Para la cooperativa y la gerencia}

\textbf{Los papeles de la ley de datos personales.} En la carpeta
\emph{legal} del proyecto quedaron preparados:
\begin{itemize}[leftmargin=1.2em]
  \item una guía que explica quién es responsable de qué ante la ley (la
  cooperativa es la responsable de los datos de su gente; vos, como operador
  del sistema, sos su «encargado»), los pasos en orden, y qué preguntarle al
  abogado;
  \item un borrador del contrato que firma cada cooperativa con vos;
  \item el formulario de consentimiento para que firme cada chofer.
\end{itemize}
Son borradores: los tiene que revisar un abogado antes de usarlos.

\subsection{Para vos, que operás el sistema}

\textbf{Saber que el servidor anda mal antes que los choferes.} El servidor
se revisa a sí mismo cada minuto (si puede leer y escribir sus datos, si le
queda espacio, si los respaldos están al día, si la hora está bien
configurada) y, si algo cambia, te puede mandar un aviso al teléfono.

\textbf{Una copia de seguridad fuera del servidor.} Cada seis horas el sistema
guarda una copia de todos sus datos. Antes esa copia quedaba en el mismo
lugar que los datos originales: si se perdía ese lugar, se perdía todo. Ahora
además sale una copia \textbf{cifrada} a un almacenamiento aparte. El
proveedor de ese almacenamiento no puede leerla. Y hay un procedimiento
probado para volver a una copia, si alguna vez hiciera falta.

\textbf{Un reinicio ya no tira las vueltas.} Cada vez que se publica una
mejora del servidor, éste se reinicia. Antes, ese reinicio le borraba a cada
combi la vuelta que venía haciendo; en hora pico eran todas. Ahora el
servidor guarda lo que tenía en curso justo antes de apagarse, y lo retoma al
volver. Las vueltas se completan como si nada hubiera pasado.

\section{La revisión a fondo}

Con tantos cambios en tan pocos días, antes de salir a la calle se hizo una
revisión completa de todo lo nuevo, hecha por cuatro revisores que no habían
escrito ese código, cada uno con una parte. Encontraron unos treinta
problemas. Cada uno se volvió a comprobar antes de tocarlo, y se corrigieron
27. Los cinco más graves:

\begin{enumerate}[leftmargin=1.5em]
  \item \textbf{Una planilla dañada podía apagar el sistema de todas las
  cooperativas.} Un archivo de Excel armado de cierta manera (incluso por
  error) hacía que el servidor se quedara sin memoria y se cayera. Ahora la
  planilla se lee aparte, con un límite: si es rara, falla esa lectura y el
  resto del sistema sigue andando.
  \item \textbf{Después de cada reinicio, un chofer podía manejar la cuenta de
  un cobrador de otra combi.} Un paso que debía hacerse una sola vez se
  repetía en cada reinicio, y le daba al chofer permisos que no tenía que
  tener. Ahora se hace una sola vez.
  \item \textbf{El archivo de datos de una persona podía traer datos de otra.}
  Si se daba de baja a alguien y después se creaba otra cuenta con el mismo
  nombre de usuario (incluso en otra cooperativa), el archivo del nuevo traía
  lo del anterior. Ahora sólo trae lo de esa cuenta.
  \item \textbf{La copia de seguridad congelaba el sistema.} Mientras se
  preparaba la copia, el sistema dejaba de atender durante varios segundos:
  ni GPS, ni emergencias, ni mapa. Ahora se prepara de a poco, sin frenar
  nada.
  \item \textbf{Dar de baja a alguien borraba lo guardado como evidencia.}
  Justo lo que no podía pasar, porque la baja suele venir después del
  incidente. Ahora la evidencia se conserva.
\end{enumerate}

Hubo otros de menor peso. Por ejemplo:
\begin{itemize}[leftmargin=1.2em]
  \item si el panel de Despacho perdía la conexión, ATENDIDO podía no llegar y
  nadie se enteraba;
  \item el motivo de una evidencia les llegaba a los teléfonos de todos los
  choferes;
  \item se podía mover un turno ya liquidado al mes actual;
  \item se podía publicar un arreglo de la aplicación sobre teléfonos que no
  estaban preparados para recibirlo.
\end{itemize}
Todos están corregidos. La lista completa está en el bloque 16 de
\emph{CAMBIOS-PARA-REVISAR.md}.

Cada arreglo quedó con su prueba automática. Antes de unir todo, se corrió la
batería completa de pruebas del sistema (más de 90) y pasó.

\section{Dónde estamos}

\begin{recuadro}{Ya está}
\begin{itemize}[leftmargin=1.2em]
  \item Todo lo de estos días está unido a la versión principal del proyecto.
  \item La lista de «lo que falta antes de salir» (los catorce puntos de la
  sección 0 de PENDIENTES) está terminada, salvo dos: el de la ley de datos
  personales, que está \emph{preparado} pero necesita trámites tuyos, y el
  de Bolivia, que quedó para después de Juliaca.
  \item Los bloqueantes para cobrarle a un cliente, de la sección 1, están
  resueltos del lado del código.
  \item El GPS con la pantalla apagada se probó 13 horas seguidas en un
  teléfono real, sin un solo corte.
\end{itemize}
\end{recuadro}

\begin{atencion}{Todavía no se probó en la calle}
Todo esto funciona en las pruebas, pero la versión 0.3.0 de la aplicación
todavía no se compiló ni se instaló en un teléfono. Lo nuevo del servidor
tampoco está todavía en Railway: llega en el próximo despliegue, ahora que
está todo unido. La primera prueba de verdad es la tuya, en tu teléfono.
\end{atencion}

\section{Lo que te toca, en orden}

Está todo en el archivo \emph{TAREAS-FUERA-DEL-CODIGO.md}, con los detalles.
Lo esencial:

\begin{enumerate}[leftmargin=1.5em]
  \item \textbf{Ver que Railway tomó la versión nueva} (o desplegarla).
  En el registro tiene que aparecer que guardó y retomó lo que tenía en
  curso.
  \item \textbf{Configurar en Railway} la zona horaria (\emph{TZ =
  America/Lima}). Sin eso, «hoy» empieza a las 7 de la tarde. Configurar
  también el correo de contacto para los datos personales y, si querés el
  aviso al teléfono, la dirección de alertas.
  \item \textbf{Configurar la copia de seguridad afuera:} crear el
  almacenamiento, poner sus datos y la frase de cifrado en Railway, y
  guardar esa frase en un lugar seguro. Sin ella las copias no se pueden
  abrir. Después, hacer un simulacro de volver a una copia.
  \item \textbf{Compilar la aplicación 0.3.0} dos veces con
  \emph{compilar-app.bat}: la de los choferes y la de tu teléfono.
  \item \textbf{Probar un arreglo sin reinstalar}: cambiar algo chico,
  publicarlo con \emph{publicar-arreglo.bat} sólo a tu teléfono, y ver que
  llega. Si llega, publicarlo a todos.
  \item \textbf{Probar en tu teléfono} lo básico: entrar, elegir la combi,
  salir a ruta, y un SOS de prueba avisándole antes a Despacho.
  \item \textbf{Ponerle un tope de gasto} a la cuenta del mapa (Geoapify).
  \item \textbf{La reunión con el abogado} por los datos personales, con la
  carpeta \emph{legal}. Después, que cada cooperativa inscriba su banco de
  datos (es gratis y en línea) y que cada chofer firme el consentimiento.
  \item \textbf{Contarle a Despacho} lo nuevo: activar la alarma al empezar el
  turno, corregir horas, cargar la planilla, guardar evidencia y bajar los
  datos de una persona.
\end{enumerate}

\section{Decisiones que te tocan}

Son cosas que el sistema deja abiertas a propósito, porque no son técnicas:

\begin{itemize}[leftmargin=1.2em]
  \item \textbf{El espacio en disco.} Ahora que las fotos se guardan un mes,
  los datos pesan más, y se guardan 28 copias locales (una semana). En el
  piloto son unos pocos gigas, alrededor de un dólar al mes. Con muchas
  cooperativas crece. La salida simple es guardar menos copias locales,
  porque la semana entera ya queda en la copia de afuera.
  \item \textbf{Cuánto se guardan las emergencias en el registro.} Hoy, un
  año. Si la cooperativa o el abogado quieren más, se cambia.
  \item \textbf{Quién puede corregir horas}: cualquier Despacho o sólo la
  gerencia. Y si al chofer se le avisa cuando le corrigen un turno.
  \item \textbf{Qué hace la cooperativa si un chofer retira su consentimiento.}
  El sistema lo anota, pero no le apaga la aplicación: esa decisión es de la
  cooperativa.
\end{itemize}

\section{Lo que viene}

\textbf{El piloto en Juliaca.} Es lo próximo, y es lo que va a decir lo que
ninguna prueba puede decir: cuánto aguanta la batería en cada marca de
teléfono, cuánto usan la pantalla los choferes, y qué les resulta difícil.

\textbf{Bolivia, en enero o febrero.} Antes de entrar a La Paz hay un día de
trabajo pendiente. Allá la hora es una más que en Perú, los números de
emergencia son otros, el formulario de consentimiento tiene que ser el de su
ley y con su carnet de identidad, y hay que generar el mapa de la ciudad. Está
anotado y se hace cuando Juliaca esté firme. Arequipa, en cambio, no necesita
nada de esto: sólo su mapa.

\vspace{2em}
\begin{center}
\small\color{gray}
Para el detalle técnico de cada cambio: \emph{CAMBIOS-PARA-REVISAR.md}.\\
Para la lista de pendientes: \emph{PENDIENTES.md}.\\
Para lo que te toca hacer: \emph{TAREAS-FUERA-DEL-CODIGO.md}.
\end{center}

\end{document}
